% Supplementary material for Section~\ref{sec:audit} (Auditing AWWL). Input
% under Appendix D (Supplementary results); starts at subsection level. Moved
% floats keep their labels (tab:replication-details is new): tab:convergence,
% fig:convergence, tab:sampler,
% tab:budget, fig:budget-sampler, fig:budget-curve, tab:deficit-budget,
% fig:deficit-metric, tab:extensions.

\subsection{Replication details}
\label{app:supp-replication}

Table~\ref{tab:replication-details} completes the FID comparison of
Table~\ref{tab:replication} with Inception Score, the one-sided tests against
the published differences and the 95\% prediction intervals for one further
paired run, under both samplers.

\IfFileExists{tables/replication_details.tex}{% Generated by scripts/make_tables.py --- do not edit by hand.
% Regenerate after every sweep; the paper reads these directly.

\begin{table}[tbp]
\centering
\footnotesize
\setlength{\tabcolsep}{3pt}
\caption{Replication details: every row of the CIFAR-10 table of the conference paper \citep{thao2026awwl} against the replication, on FID and Inception Score, epoch 199, \num{10000} samples per run. Published values are single runs (DDPM-1000, no EMA). Replicated values come from five training seeds per arm, the same weights sampled with DDPM-1000, as published, and with DDIM-100; only MSE and the two wavelet arms were sampled with DDPM-1000. Raw SNR weighting is the SNR-weighted row: the released code's unclamped weighting, which did not train stably in any seed. The Perceptual (VGG) row was not retrained (Appendix~\ref{app:setup-arms}). Differences are paired by seed, with 95\% $t$-intervals and two-sided paired $t$-tests, Holm-corrected over the two wavelet arms for DDPM-1000 (as in Table~\ref{tab:replication}) and over the 21-arm family of Table~\ref{tab:compare-fid} for DDIM-100. Bold: $p_{\mathrm{Holm}} < 0.05$ within the stated family. The one-sided test asks whether the published difference is the arm's mean effect: its $p$ is that of the alternative that the true difference is less favourable than published (larger for FID, smaller for Inception Score), uncorrected. The prediction interval asks whether it could be one run: it is the 95\% interval for the paired difference that one further run would show. A one-sided $p$ near 1 says only that the replicated difference is not less favourable than published; a replication more favourable than published (Huber on FID) shows in the prediction interval instead. In the lower three panels the Published columns repeat the published difference.}
\label{tab:replication-details}
\begin{tabular}{>{\raggedright\arraybackslash\hangindent=1em\hangafter=1}p{3.2cm}rrr@{\hspace{9pt}}rrr}
\toprule
 & \multicolumn{3}{c}{FID $\downarrow$} & \multicolumn{3}{c}{Inception Score $\uparrow$} \\
\cmidrule(lr){2-4}\cmidrule(lr){5-7}
 & Published & DDPM-1000 & DDIM-100 & Published & DDPM-1000 & DDIM-100 \\
\midrule
\multicolumn{7}{l}{\emph{Published single run; replicated mean $\pm$ sample standard deviation over five seeds}} \\
MSE & $16.68$ & $16.54 \pm 0.16$ & $18.46 \pm 0.21$ & $7.78$ & $7.834 \pm 0.055$ & $7.735 \pm 0.070$ \\
$L_1$ & $19.70$ & not run & $25.60 \pm 0.17$ & $7.69$ & not run & $7.681 \pm 0.057$ \\
Huber & $20.13$ & not run & $18.74 \pm 0.10$ & $7.62$ & not run & $7.761 \pm 0.039$ \\
Perceptual (VGG) & $20.02$ & not run & not run & $7.70$ & not run & not run \\
Raw SNR weighting & $19.83$ & not run & $325 \pm 15$ & $7.58$ & not run & $1.307 \pm 0.017$ \\
Static wavelet ($p = 0$) & $16.55$ & $16.85 \pm 0.09$ & $18.63 \pm 0.06$ & $7.81$ & $7.856 \pm 0.056$ & $7.748 \pm 0.019$ \\
AWWL & $16.62$ & $17.52 \pm 0.14$ & $18.90 \pm 0.34$ & $7.95$ & $7.790 \pm 0.069$ & $7.676 \pm 0.060$ \\
\midrule
\multicolumn{7}{l}{\emph{Difference from MSE, paired by seed, with its 95\% $t$-interval}} \\
$L_1$ & $+3.02$ & --- & $+7.14$ & $-0.09$ & --- & $-0.054$ \\
\quad 95\% CI &  & --- & $[6.79, 7.50]$ &  & --- & $[-0.181, 0.073]$ \\
\quad $p$ (raw; Holm) &  & --- & $<10^{-4}$; $\mathbf{<10^{-4}}$ &  & --- & $0.30$; $1.00$ \\
Huber & $+3.45$ & --- & $+0.28$ & $-0.16$ & --- & $+0.026$ \\
\quad 95\% CI &  & --- & $[-0.05, 0.61]$ &  & --- & $[-0.058, 0.109]$ \\
\quad $p$ (raw; Holm) &  & --- & $0.076$; $0.31$ &  & --- & $0.44$; $1.00$ \\
Raw SNR weighting & $+3.15$ & --- & $+306$ & $-0.20$ & --- & $-6.428$ \\
\quad 95\% CI &  & --- & $[287, 325]$ &  & --- & $[-6.525, -6.331]$ \\
\quad $p$ (raw; Holm) &  & --- & $<10^{-4}$; $\mathbf{<10^{-4}}$ &  & --- & $<10^{-4}$; $\mathbf{<10^{-4}}$ \\
Static wavelet ($p = 0$) & $-0.13$ & $+0.31$ & $+0.17$ & $+0.03$ & $+0.022$ & $+0.013$ \\
\quad 95\% CI &  & $[0.19, 0.43]$ & $[-0.16, 0.51]$ &  & $[-0.038, 0.083]$ & $[-0.085, 0.110]$ \\
\quad $p$ (raw; Holm) &  & $0.0022$; $\mathbf{0.0022}$ & $0.22$; $0.44$ &  & $0.36$; $0.73$ & $0.74$; $1.00$ \\
AWWL & $-0.06$ & $+0.98$ & $+0.44$ & $+0.17$ & $-0.043$ & $-0.059$ \\
\quad 95\% CI &  & $[0.80, 1.16]$ & $[-0.15, 1.03]$ &  & $[-0.191, 0.104]$ & $[-0.171, 0.054]$ \\
\quad $p$ (raw; Holm) &  & $0.0001$; $\mathbf{0.0002}$ & $0.11$; $0.32$ &  & $0.46$; $0.73$ & $0.22$; $1.00$ \\
\midrule
\multicolumn{7}{l}{\emph{One-sided test of the published difference from MSE, $p$}} \\
$L_1$ & $+3.02$ & --- & $<10^{-4}$ & $-0.09$ & --- & $0.76$ \\
Huber & $+3.45$ & --- & $1.00$ & $-0.16$ & --- & $1.00$ \\
Raw SNR weighting & $+3.15$ & --- & $<10^{-4}$ & $-0.20$ & --- & $<10^{-4}$ \\
Static wavelet ($p = 0$) & $-0.13$ & $0.0003$ & $0.032$ & $+0.03$ & $0.37$ & $0.32$ \\
AWWL & $-0.06$ & $<10^{-4}$ & $0.039$ & $+0.17$ & $0.0080$ & $0.0024$ \\
\midrule
\multicolumn{7}{l}{\emph{95\% prediction interval for the difference from MSE of one further paired run}} \\
$L_1$ & $+3.02$ & --- & $[6.28, 8.01]$ & $-0.09$ & --- & $[-0.366, 0.258]$ \\
Huber & $+3.45$ & --- & $[-0.53, 1.09]$ & $-0.16$ & --- & $[-0.179, 0.230]$ \\
Raw SNR weighting & $+3.15$ & --- & $[260, 352]$ & $-0.20$ & --- & $[-6.665, -6.190]$ \\
Static wavelet ($p = 0$) & $-0.13$ & $[0.01, 0.60]$ & $[-0.64, 0.99]$ & $+0.03$ & $[-0.126, 0.171]$ & $[-0.227, 0.252]$ \\
AWWL & $-0.06$ & $[0.54, 1.42]$ & $[-1.01, 1.89]$ & $+0.17$ & $[-0.405, 0.319]$ & $[-0.334, 0.217]$ \\
\bottomrule
\end{tabular}
\end{table}
}{%
  \begin{table}[tbp]\centering
  \fbox{\textbf{tables/replication\_details.tex not built}}
  \caption{Not built.}\label{tab:replication-details}
  \end{table}}

On Inception Score neither wavelet objective differs from MSE under either
sampler, even before correction: AWWL's paired difference is $\ddpmDISAWWL$
$[\ddpmDISAWWLLo, \ddpmDISAWWLHi]$ under DDPM-1000 (raw $p = \ddpmDISAWWLP$)
and $\dISAWWL$ $[\dISAWWLLo, \dISAWWLHi]$ under DDIM-100 (raw
$p = \dISAWWLP$), and the static wavelet arm's differences lie nearer zero.

A published single-run difference can be read as the objective's mean effect
or as the outcome of one run. A one-sided test against the published
difference, uncorrected and under the published sampler, rejects the first
reading for AWWL on both metrics (FID: $p \pPubFIDAWWLDdpm$; Inception Score:
$p = \pPubISAWWLDdpm$) and for the static wavelet arm on FID only
($p = \pPubFIDStaticDdpm$, against $p = \pPubISStaticDdpm$ for Inception
Score). The published differences were measured with DDPM-1000; compared with
DDIM-100 means, which assumes that the published effect does not depend on
the sampler, the same tests give $p = \pPubFIDAWWLDdim$ (FID) and
$p = \pPubISAWWLDdim$ (Inception Score) for AWWL.

The second reading is judged against the prediction interval. AWWL's
published Inception Score gain of $+\pubDISAWWL$ lies inside it under both
samplers, so one run can show that gain through seed variation alone. The
published FID differences lie outside it under DDPM-1000 ($\pubDFIDAWWL$
against $[\piFIDAWWLDdpmLo, \piFIDAWWLDdpmHi]$ for AWWL, $\pubDFIDStatic$
against $[\piFIDStaticDdpmLo, \piFIDStaticDdpmHi]$ for the static wavelet
arm) and inside the wider intervals of DDIM-100. Under the published sampler,
seed variation can therefore account for the published Inception Score gain
but not for the published FID differences.

\subsection{Implementation dependence of KID}
\label{app:supp-kid}

The sign of AWWL's KID difference from MSE depends on the implementation.
The library used for KID in the comparison tables
(Tables~\ref{tab:compare-other} and~\ref{tab:full}), \texttt{torch-fidelity},
gives a difference of
$+\dKIDAWWL \times 10^{-3}$ ($[\dKIDAWWLLo, \dKIDAWWLHi] \times 10^{-3}$,
uncorrected $p = \dKIDAWWLP$), whereas a
hand-written implementation with the same kernel and unbiased estimator gives
a nominally significant difference in AWWL's favour
($\dKIDhandAWWL \times 10^{-3}$,
$[\dKIDhandAWWLLo, \dKIDhandAWWLHi] \times 10^{-3}$, uncorrected
$p = \dKIDhandAWWLP$). The two differ only in choices that the definition of
KID leaves open: the Inception weights and preprocessing, a reference subset
of \num{10000} training images rather than all \num{\nRefImages}, and one
estimate over all samples rather than an average over subsets. The
\texttt{clean-fid} implementation used for the pricing test
(Section~\ref{sec:worth-headline}) gives $+\kidGapAWWL \times 10^{-3}$
($[\kidGapAWWLLo, \kidGapAWWLHi] \times 10^{-3}$, uncorrected
$p = \kidGapAWWLP$). An evaluation with the hand-written implementation alone
would have reported a gain that neither library reproduces.

\FloatBarrier
\subsection{Convergence within the 200-epoch schedule}
\label{app:supp-convergence}

The intermediate checkpoints give no evidence that AWWL reaches a given
quality with less training (Table~\ref{tab:convergence} and
Fig.~\ref{fig:convergence}). Under DDIM-100 it trails MSE at epochs 39, 79,
119 and 199 and leads only at epoch 159 (AWWL $-$ MSE $= \convDOneFiftyNine$~FID
$[\convDOneFiftyNineLo, \convDOneFiftyNineHi]$, uncorrected
$p = \convDOneFiftyNineP$, lower for \convAheadOneFiftyNine{} of the five
seeds), where seed variation is much larger than at the final checkpoint (MSE
standard deviation $\convSDMSEOneFiftyNine$~FID against
$\convSDMSEOneNinetyNine$). No checkpoint shows a significant difference in
AWWL's favour. Under the cosine learning-rate schedule these checkpoints
trace one run rather than shorter training budgets;
Section~\ref{sec:budget} and Appendix~\ref{app:supp-budget} report a longer
budget.

\IfFileExists{tables/convergence.tex}{% Generated by scripts/make_tables.py --- do not edit by hand.
% Regenerate after every sweep; the paper reads these directly.

\begin{table}[tbp]
\centering
\small
\caption{FID through training, CIFAR-10, DDIM-100, five seeds. Mean $\pm$ sample standard deviation over seeds; $\Delta$ is AWWL $-$ MSE paired by seed with its 95\% $t$-interval and two-sided paired $t$-test $p$, uncorrected. The last column counts the seeds at which AWWL has the lower FID.}
\label{tab:convergence}
\begin{tabular}{rrrrrr}
\toprule
Epoch & MSE & AWWL & $\Delta$ [95\% CI] & $p$ & AWWL ahead \\
\midrule
39 & $52.38 \pm 6.07$ & $59.91 \pm 6.13$ & $+7.53$\;$[2.97, 12.09]$ & $0.010$ & 0 of 5 \\
79 & $33.91 \pm 1.53$ & $36.91 \pm 2.61$ & $+3.01$\;$[-0.15, 6.16]$ & $0.057$ & 1 of 5 \\
119 & $24.93 \pm 1.19$ & $26.43 \pm 0.82$ & $+1.50$\;$[0.12, 2.88]$ & $0.039$ & 0 of 5 \\
159 & $22.45 \pm 3.14$ & $21.18 \pm 1.32$ & $-1.27$\;$[-4.15, 1.62]$ & $0.29$ & 4 of 5 \\
199 & $18.46 \pm 0.21$ & $18.90 \pm 0.34$ & $+0.44$\;$[-0.15, 1.03]$ & $0.11$ & 1 of 5 \\
\bottomrule
\end{tabular}
\end{table}
}{%
  \begin{table}[tbp]\centering
  \fbox{\textbf{tables/convergence.tex not built}}
  \caption{Not built.}\label{tab:convergence}
  \end{table}}

\guardedfigure{1.0}{convergence.png}{%
  FID through training for MSE, the static wavelet arm and AWWL: CIFAR-10,
  DDIM-100, \num{10000} samples, no EMA, five seeds, evaluated at epochs 39,
  79, 119, 159 and 199. (a)~Means over seeds. (b)~Mean paired
  difference from MSE with its 95\% $t$-interval, the error bars offset
  horizontally for legibility; markers are hollow as no correction is applied
  across checkpoints. Table~\ref{tab:convergence} lists the values
  for AWWL.}{fig:convergence}

\FloatBarrier
\subsection{Sampler and training budget}
\label{app:supp-budget}

Table~\ref{tab:sampler} gives the DDPM-1000 resampling of
Section~\ref{sec:sampler} with its KID differences; its FID entries repeat
those of Table~\ref{tab:replication}.

\IfFileExists{tables/sampler.tex}{% Generated by scripts/make_tables.py --- do not edit by hand.
% Regenerate after every sweep; the paper reads these directly.

\begin{table}[tbp]
\centering
\small
\setlength{\tabcolsep}{3pt}
\caption{The same trained weights under two samplers, five seeds, \num{10000} samples per run, fixed sampling seed. FID is mean $\pm$ sample standard deviation. $\Delta$ (objective $-$ MSE, paired by seed, with its 95\% $t$-interval) and $p$ refer to DDPM-1000; Holm correction is over the two arms shown, separately for FID and KID (KID from \texttt{torch-fidelity}, $\times 10^{3}$). Bold: $p_{\mathrm{Holm}} < 0.05$ within the stated family.}
\label{tab:sampler}
\begin{tabular}{lrrrrrrr}
\toprule
 & \multicolumn{2}{c}{FID} & \multicolumn{3}{c}{FID, DDPM-1000} & \multicolumn{2}{c}{KID, DDPM-1000} \\
\cmidrule(lr){2-3}\cmidrule(lr){4-6}\cmidrule(lr){7-8}
Objective & DDIM-100 & DDPM-1000 & $\Delta$ [95\% CI] & $p$ & $p_{\mathrm{Holm}}$ & $\Delta$ & $p_{\mathrm{Holm}}$ \\
\midrule
MSE & $18.46 \pm 0.21$ & $16.54 \pm 0.16$ & --- & --- & --- & --- & --- \\
Static wavelet ($p = 0$) & $18.63 \pm 0.06$ & $16.85 \pm 0.09$ & $+0.31$\;$[0.19, 0.43]$ & $0.0022$ & $\mathbf{0.0022}$ & $+0.31$ & $0.054$ \\
AWWL & $18.90 \pm 0.34$ & $17.52 \pm 0.14$ & $+0.98$\;$[0.80, 1.16]$ & $0.0001$ & $\mathbf{0.0002}$ & $+1.09$ & $\mathbf{0.0005}$ \\
\bottomrule
\end{tabular}
\end{table}
}{%
  \begin{table}[tbp]\centering
  \fbox{\textbf{tables/sampler.tex not built}}
  \caption{Not built.}\label{tab:sampler}
  \end{table}}

Table~\ref{tab:budget} and Fig.~\ref{fig:budget-sampler} set AWWL's gap to MSE
at the two budgets of Section~\ref{sec:budget} side by side. The 800-epoch
runs differ from the main matrix in three respects: the budget, an
exponential moving average of the weights, and a cosine schedule stretched to
the new length. Within each budget the two objectives are trained
identically, so the table is read row by row. At 200 epochs EMA shows no
effect that could account for the difference between the rows, although three
seeds cannot exclude one: it leaves the FID of MSE where it was ($\dMSEema$
$[\dMSEemaLo, \dMSEemaHi]$) and AWWL's gap close to its value without EMA
($+\dEMApair$, $p = \dEMApairP$, against $+\dAWWL$). Of the gaps
in Table~\ref{tab:budget}, only the 200-epoch DDPM-1000 gap is significant.

\IfFileExists{tables/budget.tex}{% Generated by scripts/make_tables.py --- do not edit by hand.
% Regenerate after every sweep; the paper reads these directly.

\begin{table}[tbp]
\centering
\small
\setlength{\tabcolsep}{3pt}
\caption{AWWL against MSE at two training budgets and two samplers; $n$ is the number of seeds. FID is mean $\pm$ sample standard deviation over seeds. Each row is one operating point, and the comparison within a row is the meaningful one: $\Delta$ is the FID gap AWWL $-$ MSE paired by seed with its 95\% $t$-interval; $p$ is the two-sided paired $t$-test, raw and Holm-corrected, for the 800-epoch rows over the two 800-epoch operating points. $^{\dagger}$The 200-epoch rows repeat tests of other families and carry those families' corrected $p$: the main family of Table~\ref{tab:compare-fid} (no EMA, DDIM-100), the one-test EMA family (EMA, DDIM-100) and the sampler family of Table~\ref{tab:replication} (DDPM-1000). Bold: $p_{\mathrm{Holm}} < 0.05$ within the stated family. The 800-epoch runs use EMA and a cosine schedule stretched to their length.}
\label{tab:budget}
\begin{tabular}{llrrrrrr}
\toprule
Budget & Sampler & $n$ & MSE & AWWL & $\Delta$ [95\% CI] & $p$ & $p_{\mathrm{Holm}}$ \\
\midrule
200 epochs, no EMA & DDIM-100 & 5 & $18.46 \pm 0.21$ & $18.90 \pm 0.34$ & $+0.44$\;$[-0.15, 1.03]$ & $0.11$ & $0.32$$^{\dagger}$ \\
200 epochs, EMA & DDIM-100 & 3 & $18.46 \pm 0.37$ & $18.95 \pm 0.20$ & $+0.49$\;$[-0.91, 1.88]$ & $0.27$ & $0.27$$^{\dagger}$ \\
800 epochs, EMA & DDIM-100 & 3 & $13.70 \pm 0.15$ & $13.85 \pm 0.10$ & $+0.15$\;$[-0.44, 0.73]$ & $0.40$ & $0.40$ \\
200 epochs, no EMA & DDPM-1000 & 5 & $16.54 \pm 0.16$ & $17.52 \pm 0.14$ & $+0.98$\;$[0.80, 1.16]$ & $0.0001$ & $\mathbf{0.0002}$$^{\dagger}$ \\
800 epochs, EMA & DDPM-1000 & 3 & $11.96 \pm 0.10$ & $12.30 \pm 0.16$ & $+0.34$\;$[-0.22, 0.91]$ & $0.12$ & $0.24$ \\
\bottomrule
\end{tabular}
\end{table}
}{%
  \begin{table}[tbp]\centering
  \fbox{\textbf{tables/budget.tex not built}}
  \caption{Not built.}\label{tab:budget}
  \end{table}}

\guardedfigure{0.72}{budget_sampler.png}{%
  AWWL's paired FID gap to MSE at the two budgets of Table~\ref{tab:budget},
  under DDIM-100 (solid line, circles) and DDPM-1000 (dashed line, squares),
  with 95\% $t$-intervals. Filled markers: $p < 0.05$ after Holm correction
  within the contrast's family (Table~\ref{tab:families}). The 800-epoch
  points rest on three seeds and differ from the 200-epoch points in EMA and
  schedule length as well as in budget.}{fig:budget-sampler}

Fig.~\ref{fig:budget-curve} follows the longer runs through training. The FID
of MSE is lowest at epoch~\longBestEpochMSE{} (\longBestFIDMSE) and rises
afterwards, so the 800-epoch endpoint lies past the optimum of this recipe;
the recipe uses no dropout, and the rise is consistent with overfitting,
which we did not test. The paired gap is positive at every evaluated
checkpoint, between $+\longGapSixNinetyNine$ and $+\longGapOneNinetyNine$,
and no checkpoint's interval excludes zero. Because the checkpoints share the
same three seed pairs, they are repeated measurements of one set of runs, not
independent replications. Averaged over the checkpoints within each seed, the
gap is $+\longGapAvg$ $[\longGapAvgLo, \longGapAvgHi]$ ($p = \longGapAvgP$).

\guardedfigure{1.0}{budget_curve.png}{%
  The 800-epoch runs at every evaluated checkpoint (epochs 199 to 799 in steps
  of 100): EMA, stretched cosine schedule, DDIM-100, three seeds.
  (a)~Seed means of FID for MSE and AWWL. (b)~Mean paired
  difference AWWL $-$ MSE, with its 95\% $t$-interval shaded. The seven
  checkpoints share the same three seed pairs.}{fig:budget-curve}

The under-training reading predicts that the baseline's high-band deficit
(Section~\ref{sec:pricing-test}) shrinks as training continues;
Table~\ref{tab:deficit-budget} tests this. Compared end against end, each run
at the annealed end of its own schedule and both with EMA, the deficit of MSE
is larger after 800 epochs than after 200, by $+\trajEndToEnd$~dB
$[\trajEndToEndLo, \trajEndToEndHi]$, while its FID improves from
\budFIDMSEShortEmaDdim{} to \budFIDMSELongDdim. Within the 800-epoch run the
deficit rises from \trajDefMSEOneNinetyNine~dB at epoch~199 to
\trajDefMSESevenNinetyNine~dB at epoch~799 (Table~\ref{tab:deficit-budget} and
Fig.~\ref{fig:deficit-metric}), but these checkpoints differ in learning-rate
phase as well as in training length, the early ones being taken at a high
learning rate, so this trajectory is descriptive only. At
epoch~\longBestEpochMSE, where FID is lowest, the deficit does not differ from
that at the end of the 200-epoch schedule (Table~\ref{tab:deficit-budget},
last column). Over this range, more training improves FID without closing the
deficit. AWWL's high-band change at the end of the longer schedule
($+\trajCorrSevenNinetyNine$~dB
$[\trajCorrSevenNinetyNineLo, \trajCorrSevenNinetyNineHi]$) has a point
estimate below its value at the end of the 200-epoch schedule with EMA
($+\trajCorrEnd$~dB $[\trajCorrEndLo, \trajCorrEndHi]$), but the two
intervals overlap; Section~\ref{sec:worth} prices the change at this budget.

\IfFileExists{tables/deficit_budget.tex}{% Generated by scripts/make_band_table.py --- do not edit by hand.
% Recomputed from the sampled images (all samples, all reference images).

\begin{table}[tbp]
\centering
\footnotesize
\setlength{\tabcolsep}{2.5pt}
\caption{High-band spectral deficit (dB; minus the signed high-band deviation of Table~\ref{tab:bands}, so larger means less high-frequency energy) and FID against training budget, all with EMA, three seeds, DDIM-100 samples. The first row is the annealed end of the 200-epoch schedule; the others are checkpoints of the 800-epoch schedule, whose early checkpoints are taken at a high learning rate. Mean $\pm$ sample standard deviation over seeds. AWWL's correction is MSE's deficit minus AWWL's, paired by seed. The last two columns give the change in MSE's deficit, paired by seed, against epoch 199 of the same run and against the end of the 200-epoch schedule (first row). Brackets are 95\% $t$-intervals.}
\label{tab:deficit-budget}
\begin{tabular}{lrrrrrr}
\toprule
 & MSE & \multicolumn{2}{c}{Deficit (dB)} & AWWL & \multicolumn{2}{c}{MSE change in deficit} \\
\cmidrule(lr){3-4}\cmidrule(lr){6-7}
Checkpoint & FID & MSE & AWWL & correction & vs.\ epoch 199 & vs.\ 200 ep. \\
\midrule
200 ep., end (199) & $18.46 \pm 0.37$ & $2.36 \pm 0.02$ & $1.95 \pm 0.09$ & $+0.41$\;$[0.14, 0.67]$ & --- & --- \\
800 ep., epoch 199 & $14.78 \pm 0.19$ & $1.85 \pm 0.03$ & $1.67 \pm 0.04$ & $+0.18$\;$[0.04, 0.33]$ & --- & $-0.51$\;$[-0.59, -0.44]$ \\
800 ep., epoch 399 & $12.01 \pm 0.17$ & $2.31 \pm 0.02$ & $1.89 \pm 0.02$ & $+0.42$\;$[0.33, 0.51]$ & $+0.46$\;$[0.31, 0.60]$ & $-0.06$\;$[-0.13, 0.02]$ \\
800 ep., end (799) & $13.70 \pm 0.15$ & $2.60 \pm 0.02$ & $2.36 \pm 0.02$ & $+0.24$\;$[0.14, 0.34]$ & $+0.75$\;$[0.68, 0.81]$ & $+0.23$\;$[0.22, 0.25]$ \\
\bottomrule
\end{tabular}
\end{table}
}{%
  \begin{table}[tbp]\centering
  \fbox{\textbf{tables/deficit\_budget.tex not built}}
  \caption{Not built.}\label{tab:deficit-budget}
  \end{table}}

\guardedfigure{0.82}{deficit_vs_metric.png}{%
  High-band deficit and FID of MSE and AWWL through the 800-epoch runs.
  (a)~High-band spectral deficit of the samples (dB below the real
  images) and (b)~FID at epochs 199, 399 and 799 (EMA, DDIM-100), mean $\pm$ one sample standard
  deviation over three seeds. The checkpoints belong to one stretched cosine
  schedule and differ in learning-rate phase as well as in training length,
  so the figure is descriptive; Table~\ref{tab:deficit-budget} compares the
  annealed endpoints of the two schedules.}{fig:deficit-metric}

\FloatBarrier
\subsection{Extensions of the weighting}
\label{app:supp-extensions}

Table~\ref{tab:extensions} gives every extension of
Section~\ref{sec:extensions} against MSE and against its parent. Apart from
GradNorm, whose global weights sum to four, no arm fixes its total weight
over $\sigma_t$. The per-sub-band, spatially adaptive and learned-basis arms
give each detail band the full $w_{\det}$, so their parent is AWWL with detail
bands summed (Eq.~\eqref{eq:summed}); the three-level arm divides $w_{\det}$
among its nine detail bands, and the remaining arms set their own weights.

The extensions do not deliver the mechanism either
(Table~\ref{tab:pricing-arms}). The per-sub-band and learned-weighting arms
leave the high band close to that of MSE, the learned basis raises it by about
as much as AWWL, and the spatially adaptive and all-axes arms lower it.
Priced by transplant (Section~\ref{sec:pricing-test}), the last two changes
are worth $\valueArmSpatial$ and $\valueArmAllAxes$~FID: each raises the FID
of the MSE samples.

The learned weighting is uncertainty weighting \citep{kendall2018multi}
conditioned on the noise level: a small network maps $\sigma_t$ to a
log-variance $\psi_b$ for each band $b$, the band is weighted by
$\tfrac{1}{2}e^{-\psi_b}$, and the total is left free (a per-band form of the
EDM2 weighting \citep{karras2024analyzing}). Each weight is then driven
towards a value inversely proportional to the loss of its band at that noise
level, which produces the steep timestep schedule reported in
Section~\ref{sec:extensions}.

Two arms failed rather than degraded, and neither is priced. GradNorm
\citep{chen2018gradnorm}, run at asymmetry $1.5$ with global weights that sum
to four, drove the LL, LH and HL weights to their positivity floor in all five
seeds, training the network through the HH band alone; it ends at
\num[round-mode=places,round-precision=1]{\fidGradNormmean}~FID with zero
recall (Table~\ref{tab:full}), a failure of the balancer at one setting
rather than evidence about balancing. The three-level arm is unstable across
seeds (\num[round-mode=places,round-precision=1]{\fidMultilevelmean}~$\pm$~%
\num[round-mode=places,round-precision=1]{\fidMultilevelsd}~FID); its
per-band averaging weights a level-$j$ detail coefficient $4^{j-1}$ times a
level-1 one (Section~\ref{sec:family-parseval}), so the instability cannot be
attributed to depth alone.

% Generated by scripts/make_tables.py --- do not edit by hand.
% Regenerate after every sweep; the paper reads these directly.

\begin{table}[tbp]
\centering
\footnotesize
\setlength{\tabcolsep}{2.5pt}
\caption{The method extensions, CIFAR-10, epoch 199, DDIM-100; $n$ is the number of seeds. No extension normalises its band weights. The per-sub-band, spatially adaptive and lifting arms keep AWWL's rational weighting with every detail band at the full $w_{\mathrm{det}}$, the summed form; the multi-level arm divides $w_{\mathrm{det}}$ among its nine detail bands; the learned-weighting, GradNorm and all-axes arms set the band weights themselves. Every extension is compared with the unnormalised AWWL with detail bands summed (second row, the parent of the summed-form arms), not with the normalised variants of Table~\ref{tab:confound}. $\Delta$ is the paired-by-seed FID difference with its 95\% $t$-interval, over the seeds both arms share; $p$ is the two-sided paired $t$-test, raw and Holm-corrected: against MSE over the 21-arm family of Table~\ref{tab:compare-fid}, against the parent over the seven extensions. Bold: $p_{\mathrm{Holm}} < 0.05$ within the stated family. Axes turned on: Per-sub-band schedules: separate schedule per detail band; Spatially adaptive: weights vary over position; Multi-level (3 levels): three-level decomposition; Learned weighting: $\sigma$-conditioned learned weights; GradNorm balancing: GradNorm sets the weights; Learned basis (lifting): learned lifting basis; All axes combined: learned weights, spatial, lifting.}
\label{tab:extensions}
\begin{tabular}{lrrrrrrr}
\toprule
 & & \multicolumn{3}{c}{Against MSE} & \multicolumn{3}{c}{Against the parent} \\
\cmidrule(lr){3-5}\cmidrule(lr){6-8}
Arm & $n$ & $\Delta$ [95\% CI] & $p$ & $p_{\mathrm{Holm}}$ & $\Delta$ [95\% CI] & $p$ & $p_{\mathrm{Holm}}$ \\
\midrule
\multicolumn{8}{l}{\emph{Reference}} \\
AWWL & 5 & $+0.44$\;$[-0.15, 1.03]$ & $0.11$ & $0.32$ &  &  &  \\
AWWL, detail bands summed & 5 & $+2.63$\;$[2.22, 3.04]$ & $<10^{-4}$ & $\mathbf{0.0008}$ & \emph{parent} &  &  \\
\midrule
Per-sub-band schedules & 5 & $+2.61$\;$[2.11, 3.11]$ & $0.0001$ & $\mathbf{0.0017}$ & $-0.02$\;$[-0.42, 0.39]$ & $0.92$ & $0.92$ \\
Spatially adaptive & 5 & $+4.65$\;$[4.10, 5.19]$ & $<10^{-4}$ & $\mathbf{0.0003}$ & $+2.02$\;$[1.70, 2.34]$ & $<10^{-4}$ & $\mathbf{0.0003}$ \\
Multi-level (3 levels) & 5 & $+88.9$\;$[50.4, 127.3]$ & $0.0030$ & $\mathbf{0.021}$ & $+86.2$\;$[47.9, 124.6]$ & $0.0034$ & $\mathbf{0.010}$ \\
Learned weighting & 5 & $+6.66$\;$[6.12, 7.20]$ & $<10^{-4}$ & $\mathbf{<10^{-4}}$ & $+4.03$\;$[3.64, 4.42]$ & $<10^{-4}$ & $\mathbf{<10^{-4}}$ \\
GradNorm balancing & 5 & $+330$\;$[296, 365]$ & $<10^{-4}$ & $\mathbf{0.0002}$ & $+328$\;$[293, 363]$ & $<10^{-4}$ & $\mathbf{<10^{-4}}$ \\
Learned basis (lifting) & 3 & $+6.13$\;$[4.99, 7.28]$ & $0.0019$ & $\mathbf{0.015}$ & $+3.36$\;$[2.37, 4.34]$ & $0.0046$ & $\mathbf{0.010}$ \\
All axes combined & 3 & $+12.5$\;$[11.5, 13.5]$ & $0.0004$ & $\mathbf{0.0039}$ & $+9.76$\;$[9.06, 10.45]$ & $0.0003$ & $\mathbf{0.0011}$ \\
\bottomrule
\end{tabular}
\end{table}


\FloatBarrier
